% Драфт подсекций 1.3.2 и 1.3.3 литературного обзора
% Секция 1.3: Методы анализа многомерной нейронной активности
% Статус: черновик для ревью

\subsection{Модели популяционной активности}\label{sec:lit_popmodels}

Методы анализа многомерной нейронной активности можно разделить на два широких класса: модели с полным наблюдением (fully observed models) и модели со скрытыми переменными (latent variable models)~\cite{Cunningham2014}. Первые описывают совместное распределение наблюдаемых нейронных откликов напрямую, вторые предполагают, что высокоразмерная активность порождается небольшим числом ненаблюдаемых факторов. Ниже кратко охарактеризованы основные представители каждого класса.

\paragraph{Модели с полным наблюдением.}
\textit{Модели максимальной энтропии} (Maximum Entropy, MaxEnt) строят наименее структурированное совместное распределение активности нейронов, удовлетворяющее заданным ограничениям --- как правило, на средние частоты разрядов и парные корреляции~\cite{SavinTkacik2017}. Модель второго порядка эквивалентна модели Изинга на графе нейронной популяции. Достоинством подхода является его минимальная предвзятость: любая дополнительная структура в данных выявляется как отклонение от MaxEnt-предсказания. Вместе с тем точный расчёт нормировочной константы требует перебора $2^N$ состояний, что делает подход вычислительно непрактичным для популяций размером более нескольких десятков нейронов.

\textit{Обобщённые линейные модели} (Generalized Linear Models, GLM) описывают вероятность спайка каждого нейрона как функцию его собственной истории, активности соседних клеток и внешних ковариат (например, параметров стимула)~\cite{Truccolo2005}. Модель хорошо масштабируется на большие популяции и допускает эффективную оценку параметров методом максимального правдоподобия. Основным ограничением является предположение о линейном характере интеграции входных сигналов: нелинейные взаимодействия между нейронами не могут быть представлены в рамках стандартной GLM.

\textit{Копула-модели} разделяют маргинальные распределения отдельных нейронов и структуру статистических зависимостей между ними~\cite{Jenison2004}. Такой подход позволяет описывать нелинейные зависимости высших порядков, сохраняя при этом гибкость в выборе модели для каждого нейрона по отдельности.

\paragraph{Модели со скрытыми переменными.}
Модели со скрытыми переменными исходят из того, что высокоразмерная нейронная активность определяется небольшим числом латентных факторов, отражающих коллективное состояние популяции. Этот подход непосредственно связан с концепцией нейронного многообразия~\cite{Gallego2017} и составляет теоретический фундамент методов, развиваемых в данной диссертации.

\textit{Статические модели} выполняют однократное снижение размерности всей записи и выявляют структуру латентного пространства без явного моделирования динамики. К этому классу относится метод jPCA~\cite{Churchland2012}, который обнаруживает вращательную динамику в моторной коре приматов при планировании движений --- латентные траектории образуют устойчивые ротации в двумерном подпространстве. Другой пример --- перекрёстно-валидированный анализ главных компонент (cvPCA)~\cite{Stringer2019}, позволивший установить степенной закон убывания дисперсии по компонентам в зрительной коре мыши: $\lambda_n \propto 1/n^\alpha$, что указывает на высокоразмерное, но структурированное кодирование стимулов.

\textit{Модели пространства состояний} (state-space models) явно описывают временну\'{ю} эволюцию латентных переменных. Наиболее известной архитектурой является LFADS (Latent Factor Analysis via Dynamical Systems)~\cite{Pandarinath2018} --- секвенциальный вариационный автоэнкодер, обученный восстанавливать пробную динамику нейронной популяции из зашумлённых одиночных проб. LFADS позволяет извлекать скрытые факторы, управляющие поведением, даже при значительном шуме и малом числе повторений. Развитием этого направления является метод VIND (Variational Inference for Nonlinear Dynamics)~\cite{ZhaoPark2017}, который допускает нелинейную динамику латентных переменных и способен находить компактные (порядка пяти измерений) латентные пространства, воспроизводящие динамику типа модели Ходжкина--Хаксли.

\textit{Модели на основе стохастических процессов} используют гауссовские процессы для совместного сглаживания и снижения размерности. Гауссов-процессный факторный анализ (GPFA)~\cite{Yu2009} оценивает гладкие латентные траектории из одиночных проб, не требуя предварительного биннинга или усреднения по повторениям. Метод GPLVM (Gaussian Process Latent Variable Model)~\cite{wu2018} обобщает этот подход, обнаруживая латентное многообразие произвольной топологии --- например, многообразие обонятельных стимулов в пириформной коре.

В целом модели со скрытыми переменными наиболее продуктивны, когда активность популяции действительно имеет низкоразмерную структуру. Их общим ограничением является то, что все источники вариативности --- как поведенчески значимые, так и шумовые --- смешиваются в едином латентном пространстве. Для содержательной интерпретации латентных факторов, как правило, необходимо явное моделирование внешних переменных, что является одной из центральных задач данной диссертации.


\subsection{Топологические подходы}\label{sec:lit_topo}

Альтернативный путь изучения геометрии нейронных многообразий предлагает топологический анализ данных (Topological Data Analysis, TDA). В отличие от методов снижения размерности, которые стремятся найти координаты в латентном пространстве, TDA выявляет качественные топологические свойства данных --- связность, наличие циклов и полостей --- через аппарат персистентных гомологий. Такие инварианты устойчивы к непрерывным деформациям и шуму, что делает топологический подход особенно ценным для нейронных данных, где точная метрическая структура зависит от выбора метода предобработки.

Ярким примером применения TDA к нейронным данным является работа Chaudhuri и соавторов~\cite{Chaudhuri2019}, в которой исследована популяционная активность переднего дорсального ядра таламуса (anterior dorsal nucleus, ADN) --- области, содержащей нейроны направления головы (head direction cells). Авторы показали, что множество состояний популяционной активности образует одномерное кольцо (топологически --- окружность $S^1$), положение на котором предсказывает текущее направление головы животного. Принципиально важно, что кольцевая структура сохраняется не только в бодрствовании, но и во время сна, когда внешний сенсорный вход отсутствует. Это свидетельствует о том, что кольцо является внутренним аттрактором нейронной цепи, а не следствием сенсорного входа, и подтверждает теоретические предсказания моделей аттракторных сетей направления головы.

Результаты TDA-анализа демонстрируют, что нейронные многообразия могут обладать нетривиальной топологией, которая не выявляется стандартными методами снижения размерности, и что эта топология может напрямую отражать структуру кодируемой информации.
